\documentclass[11pt]{article}
\usepackage[margin=1in]{geometry}
\usepackage{amsmath,amssymb,amsfonts}
\usepackage{amsthm,mathrsfs}
\usepackage{bm,bbm}
\usepackage{color}
\usepackage{xr-hyper}
\usepackage[colorlinks=true,linkcolor=blue,citecolor=blue,urlcolor=blue]{hyperref}
% labels of the main document, prefixed M-: build it first, so that its .aux exists
\externaldocument[M-]{../stratified_seir_ngm}

\theoremstyle{plain}
\newtheorem{theorem}{Theorem}[section]
\newtheorem{lemma}[theorem]{Lemma}
\newtheorem{proposition}[theorem]{Proposition}
\numberwithin{equation}{section}
\renewcommand{\thesection}{S\arabic{section}}

\title{Supplement: derivation of the next generation matrix\thanks{Parts of this manuscript were edited with the assistance of an AI tool (Claude Opus 5.5, Anthropic). The authors are solely responsible for all proofs and statements.}}
\author{Zhilan \and John \and Yuliya \and Ludmil}
\date{}

\begin{document}
\maketitle

\begin{abstract}
Supplement to ``The next generation matrix of the age- and location-stratified
SEIR model''. Equation, theorem and section numbers without a prefix S refer
to that document, whose notation is used throughout.
\end{abstract}

\section{The next generation matrix}\label{s-ngm}

For group $p$, with age group $j$, the infected states are
$(E_s,E_p,I_a,I_{ps},I_s,I_h,I_{pm},I_m)$. Linearizing the system
in~\cite{2026Feng_Day_Glasser} at the disease free equilibrium and splitting as
in~\cite{2002vdDriessche_Watmough} gives $F-V$ with $V=\operatorname{blockdiag}(V_p)$,
\begin{equation}\label{vp}
V_p=\begin{pmatrix}
\kappa&0&0&0&0&0&0&0\\
0&\kappa&0&0&0&0&0&0\\
-\kappa p_j&-\kappa q_j&\gamma_a&0&0&0&0&0\\
-\kappa(1-p_j)&0&0&\xi&0&0&0&0\\
0&0&0&-\xi&\gamma_s+\delta_s&0&0&0\\
0&0&0&0&-\delta_s&\gamma_h&0&-\delta_m\\
0&-\kappa(1-q_j)&0&0&0&0&\xi&0\\
0&0&0&0&0&0&-\xi&\gamma_m+\delta_m
\end{pmatrix},
\end{equation}
all rates evaluated at age group $j$. By the force of infection~\cite[Section~2.1]{2026Feng_Day_Glasser}
and~\eqref{M-yz}, the
only non-zero rows of $F$ are those of $E_s$ and $E_p$:
\begin{equation}\label{f-rows}
\begin{aligned}
&F_{(p,E_s),(l,\cdot)}=d_{c,p}y_pc_{pl}\psi_{I,l}\,\widehat{\bm\eta}^T,\qquad
F_{(p,E_p),(l,\cdot)}=d_{c,p}z_pc_{pl}\psi_{I,l}\,\widehat{\bm\eta}^T,\\
&\widehat{\bm\eta}=(0,0,\eta_a,\eta_{ps},1,\eta_h,\eta_{pm},\eta_m)^T .
\end{aligned}
\end{equation}
The matrix $FV^{-1}$ has order $8N$. Following~\cite{2010Diekmann_Heesterbeek_Roberts},
with $R\in\mathbb{R}^{2N\times 8N}$ selecting the components $(p,E_s)$ and
$(p,E_p)$, in the order $2p-1$, $2p$, we take
\begin{equation}\label{ngm-def}
\operatorname{NGM}=R\,FV^{-1}R^T\in\mathbb{R}^{2N\times 2N},
\end{equation}
the next generation matrix on the states at infection. It is not $FV^{-1}$,
but, as shown below, it has the same non-zero eigenvalues.

Here $\delta_s$, $\delta_m$ are the rates of hospitalization, the products
$\delta_s\gamma_{s\to h}$, $\delta_m\gamma_{m\to h}$ in~\cite{2026Feng_Day_Glasser}.
With $\phi_s=\delta_s/(\gamma_s+\delta_s)$ and $\phi_m=\delta_m/(\gamma_m+\delta_m)$, let
\begin{equation}\label{qs-qm}
\begin{aligned}
q_{s,p}&=\frac{\eta_ap_j}{\gamma_a}+\frac{\eta_{ps}(1-p_j)}{\xi}+\frac{1-p_j}{\gamma_s+\delta_s}
+\frac{\eta_h(1-p_j)\phi_s}{\gamma_h},\\
q_{m,p}&=\frac{\eta_aq_j}{\gamma_a}+\frac{\eta_{pm}(1-q_j)}{\xi}+\frac{\eta_m(1-q_j)}{\gamma_m+\delta_m}
+\frac{\eta_h(1-q_j)\phi_m}{\gamma_h},
\end{aligned}
\end{equation}
all rates evaluated at age group $j$; these are the quantities $Q_{ps}$ and
$Q_{pm}$ in~\cite[Appendix~A]{2026Feng_Day_Glasser}.

\begin{lemma}\label{q-formula}
The weights~\eqref{qs-qm} are
\begin{equation}\label{q-def}
q_{s,p}=\widehat{\bm\eta}^TV_p^{-1}\bm e_1,\qquad
q_{m,p}=\widehat{\bm\eta}^TV_p^{-1}\bm e_2 .
\end{equation}
\end{lemma}
\begin{proof}
Ordering the unknowns as $(E_s,E_p,I_a,I_{ps},I_s,I_{pm},I_m,I_h)$ makes
$V_p$ lower triangular with positive diagonal, so $V_p$ is non-singular and
$V_p\bm x=\bm e_1$ is solved in this order: $x_1=1/\kappa$, $x_2=0$;
$\gamma_ax_3=\kappa p_jx_1$ gives $x_3=p_j/\gamma_a$;
$\xi x_4=\kappa(1-p_j)x_1$ gives $x_4=(1-p_j)/\xi$;
$(\gamma_s+\delta_s)x_5=\xi x_4$ gives $x_5=(1-p_j)/(\gamma_s+\delta_s)$;
$\xi x_7=\kappa(1-q_j)x_2=0$, so $x_7=0$ and then $x_8=0$;
$\gamma_hx_6=\delta_sx_5+\delta_mx_8$ gives $x_6=(1-p_j)\phi_s/\gamma_h$.
Then $\widehat{\bm\eta}^T\bm x$ is the first line of~\eqref{qs-qm}.
For $V_p\bm x=\bm e_2$: $x_1=0$, $x_2=1/\kappa$, $x_3=q_j/\gamma_a$,
$x_4=x_5=0$, $x_7=(1-q_j)/\xi$, $x_8=(1-q_j)/(\gamma_m+\delta_m)$,
$x_6=\delta_mx_8/\gamma_h=(1-q_j)\phi_m/\gamma_h$, and
$\widehat{\bm\eta}^T\bm x$ is the second line.
\end{proof}

\begin{proposition}\label{fv}
The matrix~\eqref{ngm-def} is $\operatorname{diag}(\bm d)A_c$
of~\eqref{M-ngm}, and $FV^{-1}$ and $\operatorname{NGM}$ have the same
non-zero eigenvalues; in particular
$\rho\left(FV^{-1}\right)=\rho\left(\operatorname{NGM}\right)$.
\end{proposition}
\begin{proof}
Let $H=\operatorname{blockdiag}(\widehat{\bm\eta}^T)\in\mathbb{R}^{N\times 8N}$.
By~\eqref{f-rows}, $F=R^T\widehat F$ with
$\widehat F=D_0\operatorname{diag}(\bm d)\mathscr{P}C\operatorname{diag}(\bm\psi_I)H$.
Since $RR^T=I_{2N}$, $\operatorname{NGM}=RR^T\widehat FV^{-1}R^T=\widehat FV^{-1}R^T$.
The matrix $V^{-1}$ is block diagonal, so by Lemma~\ref{q-formula}
$HV^{-1}R^T\in\mathbb{R}^{N\times 2N}$
has in row $p$ only the entries $q_{s,p}$ in column $2p-1$ and $q_{m,p}$ in
column $2p$, i.e.\ $HV^{-1}R^T=\mathscr{P}^TQ$. With
$\operatorname{diag}(\bm\psi_I)\mathscr{P}^T=\mathscr{P}^T\Psi_I$ and the
fact that diagonal matrices commute,
$\operatorname{NGM}=D_0\operatorname{diag}(\bm d)\mathscr{P}C\mathscr{P}^T\Psi_IQ
=\operatorname{diag}(\bm d)A_c$.
Finally $FV^{-1}=R^T\left(\widehat FV^{-1}\right)$ has the same non-zero
eigenvalues as $\widehat FV^{-1}R^T=\operatorname{NGM}$, since $X_1X_2$ and
$X_2X_1$ have the same non-zero eigenvalues (Section~\ref{M-sec-reduced}).
\end{proof}

\bibliographystyle{plain}
\bibliography{../bib_mb}
\end{document}
